% =====================================================================
%  Chapter 2: AI, LLMs and multi-agent systems in education
%  (skeleton, to be written)
% =====================================================================
\chapter{Artificial Intelligence, LLMs and Multi-Agent Systems in Education}
\label{chap-sota}

\section*{Introduction}
\addcontentsline{toc}{section}{Introduction}
\todonote{Recall the conclusion of Chapter 1 and announce the plan.}

\section{Adaptive and personalized learning}
\label{sec-adaptive}
\subsection{Definitions: adaptation, personalization, individualization}
\subsection{Learner model and knowledge tracing}
\label{sec-kt}
\subsection{Intelligent tutoring systems}

\section{Context-aware learning}
\label{sec-context}
\subsection{The notion of context and its typologies}
\subsection{Learner context and task context}

\section{AI and large language models in education}
\label{sec-llm}
\subsection{From language models to LLMs}
\subsection{Educational uses: tutoring, generation, assessment}
\subsection{Retrieval-augmented generation (RAG)}
\subsection{Limits: hallucination, bias, privacy}

\section{Multi-agent systems}
\label{sec-mas}
\subsection{Agents and classical multi-agent systems}
\subsection{LLM-based multi-agent systems}
\subsection{Orchestration and collaboration patterns}

\section{Related work}
\label{sec-related-work}
\subsection{Platforms and products}
\subsection{Research prototypes}
\subsection{Comparison table}

\section{Synthesis and positioning}
\label{sec-synthesis}

\section*{Conclusion}
\addcontentsline{toc}{section}{Conclusion}
